\documentclass[10pt,a4paper]{amsart}
\usepackage[margin=19mm]{geometry}
\usepackage{amsmath,amssymb,mathtools}
\usepackage[scr=rsfs,cal=euler]{mathalfa}
\usepackage{xfrac}
\usepackage[unicode,hidelinks,bookmarks=false]{hyperref}
\newcommand{\ie}{i.e.\ }
\newcommand{\cf}{cf.\ }
\newcommand{\resp}{respectively\ }
\newcommand{\Subsection}{\subsection}
\providecommand{\nobreakdash}{\nobreak\hbox{-}}
\setlength{\emergencystretch}{2em}
\multlinegap0pt
\usepackage{enumerate}
\newcommand\NN{\mathbb{N}}
\newcommand\ZZ{\mathbb{Z}}
\newcommand\QQ{\mathbb{Q}}

\newcommand\til{\mathord\sim}

\DeclareMathOperator\Betti{Betti}
\DeclareMathOperator\Ap{Ap}
\DeclareMathOperator\supp{supp}

\newtheorem{thm}{Theorem}[section]
\newtheorem{exam}[thm]{Example}
\newtheorem{defi}[thm]{Definition}
\newtheorem{lemm}[thm]{Lemma}
\newtheorem{prop}[thm]{Proposition}
\newtheorem{rema}[thm]{Remark}
\newtheorem{nota}[thm]{Notation}
\newtheorem{conj}[thm]{Conjecture}
\newtheorem{question}[thm]{Question}
\title[Presentation cardinality in embedding dimension four]{On the cardinality of minimal presentations of numerical semigroups: original Theorem 6.3}
\author{Ceyhun Elmacioglu, Kieran Hilmer, Christopher O'Neill, Melin Okandan and Hannah Park-Kaufmann}
\date{Source: Algebraic Combinatorics 7(3) (2024), 753--771; DOI 10.5802/alco.354}
\begin{document}\maketitle
\noindent\textit{Editorial scope.} The counted unit is original Theorem 6.3 with its complete two-case authored proof. The retained uncounted core is the original Ap\'ery-set, factorization, Kunz-nilsemigroup, outer-Betti and minimal-presentation apparatus, original Theorem 2.4 and original Example 2.5, which are exactly the source-local facts the proof invokes. Only the two published vector-diagram files are not redistributed; the three Figure 6 pointers resolve by hyperlink to the published article and the proof enumerates every candidate outer Betti element explicitly in text. Printed slips, including the literal $\triangledown_B$ in the outer Betti definition, the phrase ``lies is a distinct equivalence class'' and the three-entry evaluation $\varphi(m-k^2-c,0,0)$ under the printed hypothesis $4m\geq(\eta-2)^2$, are recorded without generated corrections; every original construction, hypothesis and inequality remains.
\section{Numerical semigroups, Kunz nilsemigroups and outer Betti elements (original local core)}
Letting $m = \mathsf m(S)$, the \emph{Apery set} of $S$ is the set
$$\Ap(S) := \{n \in S: n - m \notin S\}$$
of minimal elements of $S$ within each equivalence class modulo $m$.  Since $S$ is cofinite, we~are guaranteed $|\!\Ap(S)| = m$, and that $\Ap(S)$ contains exactly one element in each equivalence class modulo $m$.  As such, we often write
$$\Ap(S) = \{a_0, a_1, \ldots, a_{m-1}\}$$
with each $a_i \equiv i \bmod m$, and view the subscripts as elements of $\ZZ_m$ (e.g. $a_{-1} = a_{m-1}$).  

Recall that a \emph{factorization} of an element $n \in S$ is an expression
$$n = z_1n_1 + \cdots + z_kn_k$$
of $n$ as a sum of atoms of $S$, which we often encode as a $k$-tuple $z = (z_1, \ldots, z_k) \in \ZZ_{\ge 0}^k$, and its \emph{length} is $z_1 + \cdots + z_k$.  
The \emph{factorization homomorphism}
$$\begin{array}{r@{}c@{}l}
\varphi_S:\ZZ_{\ge 0}^k &{}\longrightarrow{}& S \\
z &{}\longmapsto{}& z_1n_1 + \cdots + z_kn_k
\end{array}$$
is the additive semigroup homomorphism that sends each $k$-tuple $z = (z_1, \ldots, z_k)$ to the element of $S$ that $z$ is a factorization of.  Under this notation, the preimage $\varphi_S^{-1}(n) = \mathsf Z_S(n)$ is the \emph{set of factorizations} of $n \in S$.  
The \emph{kernel} of $\varphi_S$, denoted $\til = \ker\varphi_S$, relates $z \sim  z'$ whenever $\varphi_S(z) = \varphi_S(z')$, in which case we call the pair $(z,z')$ a \emph{trade} or \emph{relation}.  
The kernel is a \emph{congruence}, i.e., an equivalence relation satisfying $z + z'' \sim z' + z''$ whenever $z \sim z'$ and $z'' \in \ZZ_{\ge 0}^k$.  
A subset $\rho \subseteq \ker\varphi_S$ is called a \emph{presentation} for $S$ if the intersection of all congruences containing $\rho$ is $\ker\varphi_S$.  A presentation $\rho$ of $S$ is \emph{minimal} if no proper subset of $\rho$ is a presentation for $S$.  It is known that any two minimal presentations of $S$ have the same cardinality, which we denote $\eta(S) = |\rho|$.  
The \emph{Betti elements} of $S$ are those in the set
$$\Betti(S) := \{\varphi_S(z) : (z, z') \in \rho\},$$
where $\rho$ is any minimal presentation of $S$; the set $\Betti(S)$ is independent of the choice of $\rho$.  
The \emph{factorization graph} $\nabla_n$ of an element $n \in S$ has vertex set $\mathsf Z_S(n)$ and distinct vertices $z, z' \in \mathsf Z_S(n)$ are connected by an edge whenever $z_i > 0$ and $z_i' > 0$ for some $i$.  It is known that $n \in \Betti(S)$ if and only if $\nabla_n$ is disconnected, and in fact the number of relations $(z, z') \in \rho$ for which $n = \varphi(z)$ is one less than the number of connected components of~$\nabla_n$.  

A \emph{nilsemigroup} is a semigroup with a universally absorbing element $\infty$, called the \emph{nil}. Let $(N,+)$ be a nilsemigroup that is finite, has an identity $0 \in N$, and is \emph{partly cancellative}: $a + b = a + c \ne \infty$ implies $b = c$ for all $b, c \in N$.
Like numerical semigroups, any finite partly cancellative nilsemigroup has a unique minimal generating set.  We~write $\mathsf m(N) = |N| - 1$ for the number of non-nil elements and $\mathsf e(N)$ for one more than the number of minimal generators, which are also called \emph{atoms}.  

The \emph{Kunz nilsemigroup} $N$ of a numerical semigroup $S$ is obtained from $S/\til$, where $\til$ is the congruence that relates $a \sim b$ whenever $a = b$ or $a, b \notin \Ap(S)$ (the set $S \setminus \Ap(S)$ comprises the nil of $S/\til$), by replacing each non-nil element with its equivalence class in $\ZZ_m$.  This ensures $\mathsf e(N) = \mathsf e(S)$ and $\mathsf m(N) = \mathsf m(S)$, as the minimal generators of $N$ are the minimal generators of $S$ distinct from the multiplicity.  

A finite partly cancellative nilsemigroup $N$ may be visualized by examining the \emph{divisibility poset} $P$ of non-nil elements, wherein $b \preceq c$ when $c = a + b$ for some $a \in N$.  Partial cancellativity ensures $c$ covers $b$ when $c = a + b$ for some atom $a$, so in fact the additive structure of $N$ can be recovered from the poset structure of $P$.  If $N$ is the Kunz nilsemigroup of a numerical semigroup $S$, we call $P$ the \emph{Kunz poset} of $S$.  

Given a finite partly cancellative nilsemigroup $N$ with $\mathsf e(N) = k$, one can analogously define the factorization homomorphism $\varphi_N:\ZZ_{\ge 0}^{k-1} \to N$, with $\mathsf Z_N(p) = \varphi_N^{-1}(p)$ for each $p \in N$.  Partial cancellativity ensures $|\mathsf Z_N(p)| < \infty$ unless $p = \infty$.  
If $N$ is the Kunz nilsemigroup of a numerical semigroup $S$, then for each $a_i \in \Ap(S)$, omitting the first coordinate of each factorization in $\mathsf Z_S(a_i)$ yields $\mathsf Z_N(i)$, and $\mathsf Z_N(\infty)$ contains all remaining elements of $\ZZ_{\ge 0}^{k-1}$.  

We say $\rho \subset \ker\varphi_N$ is a \emph{minimal presentation} if it is obtained from a minimal generating set of $\ker\varphi_N$ by omitting any trade $z \sim z'$ with $\varphi_N(z) = \infty$.  Likewise, 
$$\Betti(N) := \{\varphi_N(z) : (z, z') \in \rho\}.$$
The omission of trades occurring at $\infty$ is a slight departure from the ``usual'' definition of a minimal presentation from semigroup theory, but is more natural if $N$ is the Kunz nilsemigroup of a numerical semigroup $S$.  A minimal presentation for $N$ can be obtained from a minimal presentation $\rho$ for $S$ by (i) omitting any trade $(z, z')$ with $\varphi_S(z) \notin \Ap(S)$ and (ii) omitting the leading 0 entry from both factorizations in all remaining trades.  In fact, $$\Betti(N) = \{i : a_i \in \Ap(S) \cap \Betti(S)\}.$$
There is also a partial converse.  Fix a minimal presentation $\rho$ for $S$, partitioned as $\rho = \rho' \cup \rho''$ where $(z, z') \in \rho''$ whenever $\varphi_S(z) \in \Ap(S)$.  Let $\rho'''$ be a collection of trades for $S$ obtained from a minimal presentation for $N$ by prepending a 0 entry to both factorizations in each trade.  Then $\rho' \cup \rho'''$ is also minimal presentation for $S$.  

We are now ready to define outer Betti elements.  First, the \emph{support} of a factorization $z \in \ZZ_{\ge 0}^k$ and a subset $Z \subset \ZZ_{\ge 0}^k$ are given by
$$
\supp(z) = \{i : z_i > 0\}
\qquad \text{and} \qquad
\supp(Z) = \{i : z_i' > 0 \text{ for some } z' \in Z\},
$$
respectively.  Let $\nabla_Z$ denote the graph whose vertex set is $Z$ wherein distinct vertices $z, z' \in Z$ are connected by an edge whenever $\supp(z) \cap \supp(z')$ is nonempty, and for each $i \in \supp(Z)$, define 
$$Z - e_i = \{ z - e_i :  z \in Z \text{ with } i \in \supp(z) \},$$
where $e_i$ is the $i$-th standard basis vector.  

Now, an \emph{outer Betti element} of a finite partly cancellative nilsemigroup $N$ is a subset $B \subset Z_N(\infty)$  such that
\begin{enumerate}[(i)]
\item 
for every $i \in \supp(B)$, we have $B - e_i = \mathsf Z_N(p)$ for some $p \in N \setminus \{\infty\}$, and

\item 
the graph $\triangledown_B$ is connected.
\end{enumerate}
We denote by 
$$b(N) = \text{number of outer Betti elements of } N,$$
and $\eta(N) = b(N) + |\rho|$, where $\rho$ is any minimal presentation of $N$.  

\setcounter{section}{2}\setcounter{thm}{3}
Before examining outer Betti elements in more detail, we present the following consolidation of the main results of~\cite[Section~5]{kunzfaces3} pertaining to outer Betti elements and minimal presentations.  

\begin{thm}\label{t:kunzfaces3}
If $\rho$ is a minimal presentation for the Kunz nilsemigroup $N$ of a numerical semigroup $S$, then $\rho' \cup \rho''$ is a minimal presentation for $S$, where:
\begin{enumerate}[(i)]
\item
$\rho'$ contains one trade $(z, z')$ for each outer Betti element $B$ of $N$, where $z$ is obtained from a factorization in $B$ by prepending a 0, and $z'$ is any factorization of $\varphi_S(z)$ with positive first coordinate; and

\item 
$\rho''$ is obtained from $\rho$ by prepending a 0 to both factorization of each trade.  

\end{enumerate}
In particular, $\eta(S) = \eta(N) = b(N) + |\rho|$.
\end{thm}

\setcounter{thm}{4}
\begin{exam}\label{e:outerbettis}
Any numerical semigroup $S$ with an \emph{Ap\'ery set of unique expression}, meaning that each Ap\'ery set element has exactly one factorization, has Kunz nilsemigroup $N$ whose outer Betti elements are singletons, and each contains a factorization of $\infty$ that is minimal with respect to the componentwise partial order. As such, $\eta(S) = b(N)$ by Theorem~\ref{t:kunzfaces3}.  
This is the case for the numerical semigroups $S_1$ and $S_2$ from Example~\href{https://alco.centre-mersenne.org/articles/10.5802/alco.354/}{2.1}, but not $S_3$.  

For the Kunz nilsemigroup $N_1$ of $S_1$, 
$$\mathsf Z_{N_1}(\infty) = \{z \in \ZZ_{\ge 0}^5 : z_1 + \cdots + z_5 \ge 2\},$$
so there are $\binom{6}{2}$ outer Betti elements, each containing a length 2 factorization.  
This generalizes to a known result that any max embedding dimension numerical semigroup $S$ with $\mathsf m(S) = m$ has $\eta(S) = \binom{m}{2}$, with one trade for each length 2 factorization not involving $m$.  

For the Kunz nilsemigroup $N_2$ of $S_2$, we have that $\mathsf Z_{N_2}(\infty)$ has $\eta(S) = 9$ factorizations that are minimal with respect to the component-wise partial order:\ the 5 depicted in Figure~\href{https://alco.centre-mersenne.org/articles/10.5802/alco.354/}{2}, and one for each factorization of the form $e_2 + e_i$ for $i \in [1,4]$.  Note that a finite partly cancellative nilsemigroup with 4 atoms and no other nonzero non-nil elements would have 10 outer Betti elements:\ 
\begin{itemize}
\item
one is $\{e_3 + e_4\}$, which is not an outer Betti element of $N_2$ since $6 + 7 = 5 \in N_2$;

\item
one is $\{2e_1\}$, which is not an outer Betti element of $N_2$, although $\{4e_1\}$ is an outer Betti element with identical support; and

\item
the other 8 are identical to those of $N_2$.  

\end{itemize}
These ideas are utilized in constructing the family of semigroups in Theorem~\href{https://alco.centre-mersenne.org/articles/10.5802/alco.354/}{4.2}.  

Unlike $S_1$ and $S_2$, the semigroup $S_3$ does not have an Ap\'ery set of unique expression.  Indeed, the Kunz nilsemigroup $N_3$ of $S_3$ has 3 singleton outer Betti elements, along with the outer Betti element 
$$B_3 = \{(0,2,1), (1,0,2)\}$$
which is non-singleton since it lies above $6 \in N_3$, which has 2 factorizations.  Intuitively, the factorization graph $\nabla_{70}$ has an edge between $(0,0,2,1)$ and $(0,1,0,2)$, the trade between which occurs at $46$.  This is the motivation for requirement~(ii) in the definition of outer Betti element:\ any minimal factorizations of the nil that are connected by a trade at a non-nil element cannot yield more than one trade under Theorem~\ref{t:kunzfaces3}.  
More generally, although it need not be obvious from definitions, Theorem~\ref{t:kunzfaces3} implies that prepending a 0 to any two factorizations in a given outer Betti element $B$ yields two factorizations of the same element of $S$.  
\end{exam}

\section{Minimal presentation cardinalities in embedding dimension four}
\setcounter{section}{6}\setcounter{thm}{2}
\begin{thm}\label{t:embdim4}
For any $\eta \geq 6$ and $m \in \mathbb{N}$ with $4m \geq (\eta-2)^2$, there exists a numerical semigroup $S$ with $e(S) = 4$, $\eta(S) = \eta$, and $m(S) = m$.
\end{thm}

\begin{proof}
We proceed by cases, based on the parity of $\eta$.  First, suppose $\eta = 2k + 4$ for some $k \in \ZZ_{\ge 1}$, so that $m \geq (k+1)^2$.  We claim the Kunz nilsemigroup of
$$S = \left\langle m,(k+1)m - 1,  (m-k^2-k)m +  k, (m-k^2-k)m +  k+1\right\rangle$$
is the one whose poset $P$ is depicted in Figure~\href{https://alco.centre-mersenne.org/articles/10.5802/alco.354/}{6}.  More specifically, we claim each element of $\Ap(S)$ has a unique factorization, each lying in the set 
$$A = \{(0, j, i-j, 0), (0, 0, i-j, j) : 0 \le j \le i \le k\} 
\cup \{(0, j, 0, 0) : k < j < m - k^2 - k\}.$$
As there are $2i + 1$ elements of $A$ with coordinate sum $i \le k$, it is easy to check that
$$|A| = \sum_{i=0}^{k} (2i+1) + (m - k^2 - 2k - 1) = m.$$
Additionally, each $\varphi(z)$ lies a distinct equivalence class modulo $m$ for each $z \in A$.  Indeed, if $(0,a,0,0), (0,0,b,c) \in A$ with $\varphi(0,a,0,0) \equiv \varphi(0,0,b,c) \bmod m$, then 
$$\varphi(0,0,b,c) \equiv k(b + c) + c \bmod m
\qquad \text{and} \qquad
\varphi(0,a,0,0) \equiv -a \bmod m,$$
the right hand sides of which lie in $[0, k^2 + k]$ and $[-m + k^2 + k + 1, 0]$, respectively, ensuring $a = b = c = 0$.  Meanwhile, if $(0,a,b,0), (0,0,0,c) \in A$ with $b > 0$ satisfy $\varphi(0,a,b,0) \equiv \varphi(0,0,0,c) \bmod m$, then 
$$kb - a \equiv \varphi(0,a,b,0) \equiv \varphi(0,0,0,c) \equiv kc + c \bmod m$$
necessitates $kb - a = kc + c$, and thus either (i) $b = c + 1$ and $a + c = k$, which is impossible since then $a + b = a + c + 1 = k + 1$, or (ii) $b = c + 2$ and $a = c = k$, which is impossible since then $b = k + 2$.  

It remains to show any $z \in \ZZ_{\ge 0}^4 \setminus A$ with first coordinate 0 satisfies $\varphi(z) \notin \Ap(S)$.  It suffices to assume $z$ is minimal in $\ZZ_{\ge 0}^4 \setminus A$ under the componentwise partial order, and thus corresponds to a proposed outer Betti element $B_i$ of $P$ depicted in Figure~\href{https://alco.centre-mersenne.org/articles/10.5802/alco.354/}{6}.  
Comparing equivalence classes modulo $m$, one can then readily check:\ if $i = 1$, then $z = (0,1,0,1)$ and
$$\varphi(0,1,0,1) = (m - k^2 + 1)m + k > (m - k^2 - k)m + k = \varphi(0,0,1,0);$$
if $i = 2$, then $z = (0, m - k^2-k, 0, 0)$ and 
\begin{align*}
    \varphi(0, m - k^2 - k, 0, 0) 
    &= ((m - k^2 - k)(k+1) - 1)m + k^2 + k \\
    &> k(m-k^2-k)m +  k^2 + k = \varphi(0, 0, 0, k);
\end{align*}
if $3 \le i \le k + 3$, then $z = (0, k+1-b, b, 0)$ with $b = i - 2 \in [1, k+1]$ and 
$$\varphi(0, k + 1 - b, b, 0) \ge \varphi(0, 0, b, 0) > \varphi(0, 0, 0, b - 1);$$
and if $k + 4 \le i \le 2k + 4$, then $z = (0, 0, k + 1 - c, c)$ with $c = i - k - 3 \in [1, k+1]$ and
\begin{align*}
    \varphi(0,0,k+1-c,c) 
    &= (k+1)(m-k^2-k)m + k^2 + k + c \\
    &> ((k+1)(m -k^2-k-c) - 1)m + k^2 + k + c \\
    &= (m -k^2-k-c)((k+1)m-1) = \varphi(0,m - k^2-k - c,0,0).
\end{align*}
One may now simply count the outer Betti elements of $P$ to obtain $\eta(S) = 2k + 4$.  

Next, suppose $\eta = 2k+3$ for $k \in \ZZ_{\ge 2}$, so that $m \ge k^2+k+1$.  We claim the Kunz nilsemigroup of $$S = \left\langle m, km - 1,   (m - k^2-1)m +  k, (m - k^2 -1)m +  k+1\right\rangle$$ is the one whose poset $P$ is depicted in Figure~\href{https://alco.centre-mersenne.org/articles/10.5802/alco.354/}{6}. More specifically, we claim each element of $\Ap(S)$ has a unique factorization, each lying in the set 
\begin{align*}
A &= \{(0, j, i-j, 0) : 0 \le j \le i \le k\}
\cup \{(0, 0, i-j, j): 1 \le j \le i \le k-1\} \\
& \qquad \cup \{(0, j, 0, 0) : k < j \le m - k^2 -1\}.
\end{align*}
As there are $2i+1$ elements of $A$ with coordinate sum $i \leq k-1$, we see
$$|A| = \sum_{i=0}^{k-1} (2i+1) + (k + 1) + (m-k^2-k - 1) = m.$$ Additionally, we see that $\varphi(z)$ lies is a distinct equivalence class modulo $m$ for each $z \in A$. Indeed, if $(0,a,0,0), (0,0,b,c) \in A$ with $\varphi(0,a,0,0) \equiv \varphi(0,0,b,c) \bmod m$, then 
$$\varphi(0,0,b,c) \equiv k(b + c) + c \bmod m
\qquad \text{and} \qquad
\varphi(0,a,0,0) \equiv -a \bmod m,$$
the right hand sides of which lie in $[0, k^2]$ and $[-m + k^2 + 1, 0]$, respectively, thereby ensuring $a = b = c = 0$. Meanwhile, if $(0,a,b,0), (0,0,0,c) \in A$ with $b > 0$ satisfy $\varphi(0,a,b,0) \equiv \varphi(0,0,0,c) \bmod m$, then 
$$kb - a \equiv \varphi(0,a,b,0) \equiv \varphi(0,0,0,c) \equiv kc + c \bmod m$$
necessitates $kb - a = kc + c$, and thus either (i) $b = c + 1$ and $a + c = k$, which is impossible since then $a + b = a + c + 1 = k + 1$, or (ii) $b = c + 2$ and $a = c = k$, which is impossible since then $b = k + 2$.  

Proceeding as before, it remains to show that for each proposed outer Betti element $B_i = \{z\}$ of $P$, we have $\varphi(z) \notin \Ap(S)$.  
If $i = 1$, then $z = (0,1,0,1)$ and
$$\varphi(0,1,0,1) = (m - k^2 + k + 1)m + k > (m - k^2-1)m +  k = \varphi(0,0,1,0);$$
if $i = 2$, then $z = (0,m - k^2,0,0)$ and 
\begin{align*}
    \varphi(0, m - k^2, 0, 0) 
    &= (km - k^3 - 1)m + k^2 
    > k(m-k^2-1)m +  k^2 = \varphi(0, 0, 0, k);
\end{align*}
if $3 \le i \le k + 2$, then $z = (0, k+1-b, b, 0)$ with $b = i - 2 \in [1, k+1]$ and 
$$\varphi(0, k + 1 - b, b, 0) \ge \varphi(0, 0, b, 0) > \varphi(0, 0, 0, b - 1);$$
if $i = k+3$, then $z = (0,0,k+1,0)$ and 
\begin{align*}
    \varphi(0,0, k+1,0) &= (k+1)(m - k^2 - 1)m + k^2 + k\\
    &> (km - k^3 - k^2 - 1)m + k^2 + k = \varphi(0,m-k^2-k,0,0);
\end{align*}
and if $k+4 \le i \le 2k+3$, then $z = (0,0,k-c,c)$ with $c = i - k - 3 \in [1,k]$ and 
\begin{align*}
    \varphi(0,0,k-c,c) &= (m-k^2+k)km + k^2 + c\\
    &> (km - k^3 - ck - 1)m + k^2 + c\\
    &= (m - k^2 - c)(km - 1) = \varphi(m - k^2 - c,0,0).
\end{align*}
Counting the outer Betti elements of $P$ yields $\eta(S) = 2k + 3$, as desired.
\end{proof}

\begin{thebibliography}{99}
\bibitem[8]{kunzfaces3} Tara Gomes, Christopher O'Neill, and Eduardo Torres Davila, ``Numerical semigroups, polyhedra, and posets III: Minimal presentations and face dimension'', \emph{Electron. J. Combin.}\textbf{30}(2023),no.2, article no.2.57 (25pages), https://doi.org/10.37236/10380.
\end{thebibliography}
\end{document}
